\section{方法}

本文在\textbf{统一实现}（同一执行器、同一通过判定与统计口径、同一编排入口）下实现\textbf{预算分档编排}：总体形态见图~\ref{fig:overall_framework}，按题执行的端到端流水线及四档差异见第\ref{sec:overview-flow}节。其后先给出问题边界与成本记账（第\ref{sec:formal-objective}节），再分节展开难度分流、失败救援与级联等模块，用以解释各档内部的计算投向。

本文的方法设计围绕一个基本判断展开：代码生成中的额外测试时计算，只有在能够显著缩小后续搜索空间时才具有实际价值\cite{chen2024selfdebug,shinn2023reflexion,gou2024critic,li2025sstar}。因此，系统并未把预算简单地理解为「增加若干次调用」，而是将其分配到三个彼此衔接的环节中：首先判断哪些样本值得进入重路径，其次在第一次失败后提取能够约束后续生成的局部错误信息，最后在更高预算下将计算扩展到候选探索与筛选层面。如图~\ref{fig:overall_framework} 所示，本文方法的核心不在于无差别地增加阶段，而在于让不同预算档位对应不同的信息利用方式，使各档具有清晰的功能分工，而不是演化为若干边界含混的工程开关。

\begin{figure*}[t]
\centering
\includegraphics[width=0.98\textwidth]{gpt-fig2-overall-framework.png}
\caption{预算分档编排整体框架}
\label{fig:overall_framework}
\end{figure*}


\subsection{总体流程与预算分档}
\label{sec:overview-flow}

对每道任务 $x_i$，系统在实现上遵循同一\textbf{外层闭环}：读入题面及（对模型可见的）测试接口信息 $\rightarrow$ 语言模型生成候选程序 $y_i$ $\rightarrow$ 在沙箱中执行单测并得到通过/失败 $\rightarrow$ \textbf{若失败}，则仅在当前\textbf{预算档位}所允许的范围内追加后续调用（修补、重链、救援或级联），直至通过或达到该档策略的上限。通过率、平均调用次数与总 token 均在同一统计管道中汇总，从而保证四档结果可比。

分档由单一入口控制：实现中将离散档位编码与仓库内模式名一一对应，避免「模式名」与「预算档位」两套开关并存带来的比较歧义。映射见表~\ref{tab:tier_mapping}；完整命令行见仓库复现说明。

各档在行为上可概括为：
\begin{itemize}
\item \textbf{单步}：一次生成即提交判定；该档不把失败样本继续送入本文后续多轮编排分支，用作轻量基线与残差起点。
\item \textbf{单步+修补}：在单次生成失败后，以执行反馈为约束进行有限轮次的修补（轻路径；对应 \texttt{single\_step\_repair} 形态）。
\item \textbf{分流+重链与救援}：先依据结构难度分数 $s_i$ 与阈值 $\tau$ 划分 easy/hard；easy 侧保持相对轻量，hard 侧进入更重主链（可含 spec/cot 等开关），并在增强修复阶段按需触发候选救援与 refine 救援；失败信息经压缩后进入修补/替换等闭环（见第\ref{sec:rescue}节）。
\item \textbf{级联}：在更高预算下扩展多候选生成、执行筛选与 refine；主实验采用级联默认基线，将预算集中于主级联链路（见第\ref{sec:cascade}节）。
\end{itemize}

同一外层闭环在四档下复用相同的执行与判定逻辑；各档差异主要体现在失败后的\textbf{追加调用是否发生、以何种提示形态发生、以及是否引入多候选级联}。下文第\ref{sec:formal-objective}节给出问题边界与成本记账，其后按模块说明分流、救援与级联。

\subsection{问题定义与成本记账原则}
\label{sec:formal-objective}

需要强调的是，上述编排逻辑与具体基准数据集\textbf{正交}：同一套执行器、同一套分档入口与同一统计口径同时作用于 HumanEval 与 MBPP；差异只体现在任务分布与难度天花板上，而这正是双基准并列评测的意义——观察同一编排是否在“易饱和”与“更难、更杂”的两类代码生成场景下给出一致的趋势与可解释的成本结构。

设数据集为 $\mathcal{D}=\{x_i\}_{i=1}^{N}$，其中 $x_i$ 为一道代码生成任务。系统在给定预算档位 $b$ 下生成候选 $y_i$，经执行器得到 $z_i=E(y_i)\in\{0,1\}$。本文关心的不是写出抽象的约束优化闭式解，而是\textbf{在可测的记账口径下}比较各档：对每题记录是否通过、调用次数与总 token；主文以\textbf{总 token}作为 API 成本主代理（与第~\ref{sec:cost_rho} 节 $\rho$ 一致），调用次数作为辅助指标。\textbf{单次生成}提供最低档基线，并决定后续各档所面对的\textbf{残差失败分布}——后续阶段并非替代单步，而是对残差样本做有预算上限的追加计算。

本文不在正文引入带数值权重 $\alpha,\beta$ 的联合成本泛函；若需同时按次计费与按 token 计费，可将实验中的 $\rho$ 定义改为「在固定 $\alpha/\beta$ 下的边际收益」，作为部署扩展而非本文主表默认。

\subsection{难度感知路由}
\label{sec:routing}

难度感知路由的作用，是在进入重路径之前完成一次粗粒度预算筛分。系统依据题面与测试可见结构特征，对样本进行 easy/hard 划分，使明显简单样本尽量停留在轻路径，而将更可能需要额外推理或修补的样本送入重路径。这里追求的不是完美预测器，而是避免把同等预算平均花在所有样本上。

记样本 $x_i$ 的结构难度分数为 $s_i$。实现上由 \texttt{\_difficulty\_score\_v2} 给出可分解启发式：五类子分数分别刻画题面复杂度、签名/接口复杂度、测试用例复杂度、边界条件密度与题面歧义程度；另叠加\textbf{有界体量项} $e_i$（由题面/测试长度、行数、断言行数、\texttt{candidate} 调用次数等经对数与线性缩放后相加并\textbf{上限截断}，避免总分分布漂移）。离线统计表与主程序使用同一分解，便于附录核对。
\begin{equation}
\label{eq:difficulty}
\begin{aligned}
s_i={}& p_i+\sigma_i+\tau_i+b_i+a_i + e_i,
\end{aligned}
\end{equation}
其中 $p_i,\sigma_i,\tau_i,b_i,a_i$ 分别对应题面、签名、测试、边界与歧义五类子项，$e_i$ 为体量增量。实现上，五类子分项各自在代码中经\textbf{上界裁剪}后相加得到基础分，再叠加有界体量项 $e_i$；各子项内部由若干启发式规则与\textbf{固定正权重}线性组合（例如题面长度分段、关键词命中按次加权并封顶、\texttt{assert} 行计数分段、\texttt{candidate} 形参与边界词表加权等），总分为\textbf{确定性函数}而非学习到的预测器。相对权重体现为：测试结构与边界信号在总分中通常占比较高，题面歧义与体量项起校准尺度作用；子分项与结构分档、运行标签的统计关系见附录表~\ref{tab:appendix_weight_evidence}。本文将路由解释为\textbf{预算分配规则}：用 $s_i$ 与 $\tau$ 控制 hard 覆盖率，从而在精度与 token 之间形成可调前沿（阈值消融见实验节）。

路由决策写为
\begin{equation}
r_i=\mathbb{I}(s_i\ge \tau_{\mathrm{eff}}),
\end{equation}
其中 $\tau_{\mathrm{eff}}$ 与配置中的名义结构阈值一致（当前分流+重链与救援档默认取 $1.6$）；实现仅统计分数分布（最小/中位/最大）供分析，\textbf{不对}阈值做“分布地板抬升”。$r_i=1$ 表示样本进入 hard 主链。在 \texttt{tiered} 等模式下还可叠加 LLM 评审等信号，本文主表仍聚焦于上述结构分驱动的分档端到端效果。

\textbf{Multiset 随机路由对照（实现口径）。}为分离「hard 槽位覆盖率」与「槽位与题面结构的对应关系」，仓库提供 \texttt{--difficulty\_assignment=random\_permute\_matched}：先按式~\eqref{eq:difficulty} 与 $\tau_{\mathrm{eff}}$ 得到本批任务上的 easy/hard 标签 multiset（各题 structural 判定），再在固定随机种子下对该标签序列做\textbf{均匀随机置换}，使每题被赋予 easy 或 hard 的\textbf{个数}与 structural 完全一致，但题—槽位对应被打乱；随后仍执行同一套分流+重链与救援协议。该对照\textbf{不是}「任意 token 预算下的均匀随机路由」，而是\textbf{在固定 hard 覆盖率 multiset 下检验结构对齐是否带来额外收益}，与 LR 离线证据回答的问题正交互补。

\textbf{可复现性与随机性。}随机种子固定任务顺序与部分随机模块，但生成调用仍采用非零温度，同配置重复运行可能出现个别样本通过与否或总 token 的波动；表~\ref{tab:main_results} 报告的是三 seed 的\textbf{均值与样本标准差}。部分补充表为控制算力仍为单次 seed，正文在对应表下注明。

\textbf{关于体量项与附录统计。}$e_i$ 的作用是让总分与题面、测试规模保持弱对齐，便于路由分数在不同题型下保持可解释尺度。附录中的相关性与逻辑回归结果用于描述这一结构信号与运行标签的关系；主结论仍由同协议端到端结果给出，即通过率、调用次数与 token 的联合变化。

这一阈值设定使语义保持稳定：调高阈值使更多样本停留在 easy 链，调低则更多进入 hard 链；后文阈值消融亦建立在“直接阈值、直接分流”的实现上。

\subsection{失败后救援链}
\label{sec:rescue}

难度感知只能决定哪些样本进入重路径，但并不能保证额外预算转化为有效收益。为此，本文在\textbf{分流+重链与救援}档的重路径中引入失败后救援链。其思路不是简单重试，而是将执行失败转化为后续生成可利用的局部约束\cite{chen2024selfdebug,shinn2023reflexion,gou2024critic}。系统先执行候选程序，再从运行结果中提炼与修补直接相关的错误信息，并回写到下一轮生成中。

这一设计与普通重复采样的根本区别在于，后续调用不再面对原始问题，而是面对一个已经被错误信息裁剪过的更窄搜索空间。对代码生成任务而言，第一次失败通常已经暴露出相当具体的局部结构，例如参数个数不匹配、断言不满足、边界条件处理错误，或实现主体基本正确但中间步骤存在偏差。若这些信息能够被压缩为紧凑且具有约束性的反馈，则第二次生成更可能表现为局部修补，而不是从头重写。

基于这一考虑，本文将救援链组织为“生成、执行、回写错误、再生成”的短闭环，并强调对错误信息的结构化整理，而不是直接复制全部运行日志。实现上，系统在 repair 与 replace 之间按错误类型分流：当失败形态呈现“局部可修复”特征时，如断言差异、边界条件遗漏、变量更新错误，优先走 repair；当失败形态表现为“整体方案偏离”时，如函数骨架与测试意图不匹配、核心思路错误或重复修补仍不收敛，转入 replace。该分流的目的，是把额外调用优先用于最可能产生增量的路径，避免把全局重写与局部修补混在同一提示中。

\textbf{分流机制与输入。}\texttt{repair}/\texttt{replace} 标签由\textbf{规则函数}对执行器返回的 \texttt{check\_correctness} 结果字符串做关键词判定得到（非额外模型调用）：例如未闭合字符串、\texttt{InvalidSyntax} 等归入 \texttt{replace}；缺参、参数个数不匹配、\texttt{NameError}/\texttt{IndentationError} 等可归入 \texttt{repair}；其余默认 \texttt{replace}。在 direct 与 spec 双路径都失败时，实现还比较两条路径错误标签是否一致，以决定后续修补轮数为 1 或 2（一致视为“易定位”，不一致视为“分歧、难定位”）。第三档进一步提供两级救援——候选救援与 refine 救援，可分别开关。其中 \texttt{hard\_spec} 的角色是为 hard 路径提供\textbf{结构化约束提示}，即接口、边界与测试意图，使第一次重链生成先收窄解空间再进入后续救援；因此它属于主链约束模块而非后处理修饰模块。记第一次失败反馈为 $e_i^{(1)}$，压缩后的回注包为
\begin{equation}
\phi(e_i^{(1)})=\operatorname{Concat}\bigl(
\textit{错误类型},\,
\textit{截断摘要},\,
\textit{规则片段},\,
\textit{定向修补提示}
\bigr),
\end{equation}
读者可将四项理解为：归类信号、可放入提示的短错误摘要、从历史失败归纳的修补规则、以及指向具体位置的修补提示（实现中由统一格式化函数打包）。因此第二轮生成上下文可表示为
\begin{equation}
\tilde{x}_i^{(2)}=\operatorname{Concat}\bigl(x_i,\phi(e_i^{(1)})\bigr).
\end{equation}
相应地，第二轮候选程序满足
\begin{equation}
y_i^{(2)} \sim p_{\theta}\bigl(\cdot \mid \tilde{x}_i^{(2)}\bigr).
\end{equation}

这一写法强调的不是模型形式本身，而是搜索空间在执行反馈作用下发生了变化。原始 traceback 往往冗长且包含大量与修补无直接关系的细节，未经处理地回注模型容易重新放大搜索空间。该主张在表~\ref{tab:e2_feedback_format_5seed} 的 E2 对照中获得独立检验：在仅替换回注形态、其余流程不变时，\texttt{compressed} 与 \texttt{raw\_traceback} 的五 seed 均值通过率差异很小，而 \texttt{compressed} 的平均总 token 更低。因而，失败后救援的关键不在于是否追加了一次调用，而在于是否将追加预算转化为更强的局部约束。

从整体方法上看，分流+重链与救援档之所以应当被视作一条独立路径，正是因为它改变了额外预算的作用位置。相比于仅在题目层面继续扩展推理，失败后救援把预算进一步下沉到了错误层面的局部修补之中。这种变化虽然在流程图上只体现为少量附加步骤，但对系统搜索行为的影响是实质性的。

\begin{figure}[t]
\centering
\includegraphics[width=0.98\linewidth]{gpt-fig4-rescue-loop.png}
\caption{失败反馈救援闭环}
\label{fig:rescue_loop}
\end{figure}


\subsection{高预算级联工作流}
\label{sec:cascade}

级联模块（\texttt{workflow\_cascade}）的设计目标不是继续加深单样本修补链，而是把新增计算用于\textbf{候选空间覆盖}。其架构可分为三层：\textbf{候选扩展层}（生成多候选并拉开解空间）、\textbf{执行筛选层}（在统一执行器下按测试结果过滤）、\textbf{后处理层}（仅对进入候选集的样本做有限 refine）。该设计把“生成--验证--再组织”写成显式流水线，避免把多阶段流程退化为无差别重试\cite{li2022alphacode,li2025sstar}。

实现上，级联主链按以下顺序执行：\textbf{(1) direct 基线候选}给出低成本起点；\textbf{(2) \texttt{extra} 候选扩展}在失败或低置信样本上追加候选集合；\textbf{(3) 执行筛选}按统一判定口径选择最优可行候选；\textbf{(4) refine/fuse}仅在候选池已形成可行解时触发一次性微调。该顺序的目的，是保证主要预算先用于“扩展可行解覆盖率”，而非提前消耗在后处理上。

\textbf{候选数量、可行解与 refine。}\textbf{可行解}指在给定任务上于统一执行器下通过官方单元测试的候选实现。\texttt{extra} 阶段每题最多追加的生成次数由全局上限 \texttt{max\_repair\_rounds} 与路由表共同约束：在可学习路由的自适应档位绑定中，保守/平衡/进取档分别对应最多约 1/2/3 个额外生成，并对执行器实际执行的 top-$k$ 做同档绑定（例如进取档最多执行前 3 个生成候选）；启发式路由下仍受 CLI 默认 \texttt{extra\_exec\_topk} 等下限约束，实际执行数为 $\min(\text{生成数},\text{top-}k)$。\textbf{refine\_once} 仅在开启相应 CLI 开关且路由允许时触发，用于在已形成可行候选池后对最优候选做一次微调；保守档可在路由表中关闭 refine 以省 token。

与第\ref{sec:rescue}节的失败后救援链相比，级联模块的职责边界是清晰分开的：救援链解决“已知失败样本如何被局部约束修复”，级联模块解决“在不确定解空间中如何提高命中率”。因此，级联并不是救援链的机械叠加，而是把计算投向从局部修补切换到候选覆盖与筛选组织；两者解决的是不同层级的问题。

实现上，级联模块遵循“先筛选、后精修”的触发顺序：只有候选在执行筛选后保留，才进入 refine；当候选扩展未带来有效可行解时，流程在该层停止而不继续放大后处理调用。结合实验节级联开关消融，\texttt{extra}（候选扩展）提供主增益，refine 相关开关主要负责精度--成本微调。方法上应将级联理解为“用候选覆盖换取上限”，而非对中等预算链路的逐层替代。

\subsection{预算分档与模块映射}

综合而言，\textbf{预算分档编排}将\textbf{从单步到级联}组织为四个可对照层级：\textbf{单步}提供轻量起点；\textbf{单步+修补}把计算投向可压缩的执行反馈；\textbf{分流+重链与救援}在难度分流下决定重链覆盖范围，并由失败后救援把追加调用转化为局部约束；\textbf{级联}则把计算主要投向候选扩展、执行筛选与 refine。各档差异对应的是\textbf{预算投向与信息利用方式}，而非孤立工程开关。

\begin{table*}[htbp]
\centering
\small
\caption{语义档位与实现模式映射（正文口径）}
\label{tab:tier_mapping}
\begin{tabularx}{\textwidth}{p{2.6cm} p{3.5cm} X}
\toprule
语义档位 & 实现模式 & 主要计算投向 \\
\midrule
单步 & \texttt{single\_step} & 单次生成与执行判定 \\
单步+修补 & \texttt{single\_step\_repair} & 执行反馈约束下的轻量修补 \\
分流+重链与救援 & \texttt{difficulty\_aware\_plus} & 难度分流、hard 主链与失败后救援 \\
级联 & \texttt{workflow\_cascade} & 多候选扩展、执行筛选与 refine \\
\bottomrule
\end{tabularx}
\vspace{0.5ex}
{\footnotesize\textbf{注：}第三档在双基准上共用同一实现 \texttt{run\_mbpp\_difficulty\_aware\_plus}；命令行预算档位别名为 \texttt{difficulty\_aware\_repair}，仅用于复现，不改变协议定义。第一档日志别名 \texttt{base\_direct} 与 \texttt{single\_step} 同义。}
\end{table*}

这一点也是本文实验比较成立的前提：读者在阅读表~\ref{tab:main_results} 时，应始终把每一行理解为\textbf{同一实现下某一档位的端到端形态}，从而把性能差异解释为分档链条与任务分布的相互作用，而不是孤立模块的偶然波动。